% from ega2-1-fr.tex
\oldpage[II]{65}

    le faisceau gradué quasi-cohérent d'idéaux de $\mathcal{S}$, image
    réciproque de $\Gamma_*(\mathcal{I})$ par l'homomorphisme canonique
    $\alpha:\mathcal{S}\to\Gamma_*(\mathcal{O}_X)$
    (\hyperref[II.3.3.2-fr]{3.3.2}), et posons
    $\mathcal{S}'=\mathcal{S}/\mathcal{J}$. Alors $X'$ est le sous-préschéma
    associé (\hyperref[I.4.2.1-fr]{I, 4.2.1}) à l'immersion fermée
    $\operatorname{Proj}(\mathcal{S}')\to X$ correspondant à l'homomorphisme
    canonique $\mathcal{S}\to\mathcal{S}'$ de $\mathcal{O}_Y$-Algèbres
    graduées.
\end{enumerate}
\end{proposition}

\begin{proof}
\begin{enumerate}
  \item[\rm(i)] On peut supposer $\varphi$ surjectif
    (\hyperref[II.3.6.1-fr]{3.6.1}). Alors, pour tout ouvert affine $U$ de
    $Y$, $\Gamma(U,\mathcal{S})\to\Gamma(U,\mathcal{S}')$ est surjectif
    (\hyperref[I.1.3.9-fr]{I, 1.3.9}), donc
    (\hyperref[II.3.5.1-fr]{3.5.1}) on a $G(\varphi)=X$. On est
    immédiatement ramené à démontrer la proposition lorsque $Y$ est affine,
    et alors elle découle de (\hyperref[II.2.9.2-fr]{2.9.2, (i)}).
  \item[\rm(ii)] On est ramené à prouver que l'homomorphisme
    $\widetilde{\mathcal{J}}\to\mathcal{O}_X$ déduit de l'injection canonique
    $\mathcal{J}\to\mathcal{S}$ est un isomorphisme de
    $\widetilde{\mathcal{J}}$ sur $\mathcal{I}$~; comme la question est locale
    sur $Y$, on peut supposer $Y$ affine d'anneau $A$, ce qui entraîne
    $\mathcal{S}=\widetilde{S}$, où $S$ est une $A$-algèbre graduée engendrée
    par $S_1$, $S_1$ étant de type fini sur $A$. Il suffit alors d'appliquer
    (\hyperref[II.2.9.2-fr]{2.9.2, (ii)}).
\end{enumerate}
\end{proof}

\begin{corollary}[3.6.3]
\label{II.3.6.3-fr}
Sous les conditions de (\hyperref[II.3.6.2-fr]{3.6.2, (i)}), et en supposant
de plus $\mathcal{S}$ engendrée par $\mathcal{S}_1$,
$\Phi^*(\mathcal{O}_X(n))$ s'identifie canoniquement à
$\mathcal{O}_{X'}(n)$ pour tout $n\in\mathbf{Z}$.
\end{corollary}

\begin{proof}
On a défini un tel isomorphisme canonique lorsque $Y$ est affine
(\hyperref[II.2.9.3-fr]{2.9.3})~; dans le cas général, il suffit de vérifier
que les isomorphismes ainsi définis pour chaque ouvert affine $U$ de $Y$ sont
compatibles avec le passage de $U$ à un ouvert affine $U'\subset U$, ce qui
est immédiat.
\end{proof}

\begin{corollary}[3.6.4]
\label{II.3.6.4-fr}
Soient $Y$ un préschéma, $\mathcal{S}$ une $\mathcal{O}_Y$-Algèbre graduée
quasi-cohérente engendrée par $\mathcal{S}_1$, $\mathcal{M}$ un
$\mathcal{O}_Y$-Module quasi-cohérent, $u$ un $\mathcal{O}_Y$-homomorphisme
surjectif $\mathcal{M}\to\mathcal{S}_1$,
$\overline{u}:\mathbf{S}_{\mathcal{O}_Y}(\mathcal{M})\to\mathcal{S}$
l'homomorphisme de $\mathcal{O}_Y$-Algèbres graduées prolongeant $u$
(\hyperref[II.1.7.4-fr]{1.7.4}). Le morphisme correspondant à
$\overline{u}$ est alors une immersion fermée de
$\operatorname{Proj}(\mathcal{S})$ dans
$\operatorname{Proj}(\mathbf{S}_{\mathcal{O}_Y}(\mathcal{M}))$.
\end{corollary}

\begin{proof}
En effet, $\overline{u}$ est surjectif par hypothèse et on applique
(\hyperref[II.3.6.1-fr]{3.6.1, (i)}).
\end{proof}

\subsection{Morphismes d'un préschéma dans un spectre homogène.}
\label{subsection:II.3.7-fr}

\begin{env}[3.7.1]
\label{II.3.7.1-fr}
Soient $q:X\to Y$ un morphisme de préschémas, $\mathcal{L}$ un
$\mathcal{O}_X$-Module inversible, $\mathcal{S}$ une $\mathcal{O}_Y$-Algèbre
graduée quasi-cohérente à degrés positifs~; $q^*(\mathcal{S})$ est alors une
$\mathcal{O}_X$-Algèbre graduée quasi-cohérente à degrés positifs.
Considérons la $\mathcal{O}_X$-Algèbre graduée quasi-cohérente
$\mathcal{S}'=\bigoplus_{n\geq0}\mathcal{L}^{\otimes n}$ et supposons donné
un $\mathcal{O}_X$-homomorphisme d'Algèbres graduées
\[
  \psi:q^*(\mathcal{S})\longrightarrow
  \mathcal{S}'=\bigoplus_{n\geq0}\mathcal{L}^{\otimes n}
\]
ce qui revient d'ailleurs à se donner un $q$-morphisme d'Algèbres graduées
\[
  \psi^\flat:\mathcal{S}\longrightarrow q_*(\mathcal{S}').
\]
On sait que $\operatorname{Proj}(\mathcal{S}')$ s'identifie canoniquement à
$X$ (\hyperref[II.3.1.7-fr]{3.1.7} et
\hyperref[II.3.1.8-fr]{3.1.8, (iii)})~; on déduit canoniquement de $\psi$ un
ensemble ouvert $G(\psi)$ de $X$ et un $Y$-morphisme
\[
\label{II.3.7.1.1-fr}
  r_{\mathcal{L},\psi}:G(\psi)\longrightarrow
  \operatorname{Proj}(\mathcal{S})=P
\tag{3.7.1.1}
\]

